\documentclass[10pt]{article}
\usepackage{tmlr}
% Anonymous by default. For the de-anonymised preprint use:
%   \usepackage[preprint]{tmlr}
% If accepted:
%   \usepackage[accepted]{tmlr}

\usepackage{amsmath,amssymb}
\usepackage{graphicx}
\usepackage{booktabs}
\usepackage{lmodern}
\usepackage[expansion=false]{microtype}
\usepackage[table]{xcolor}
\usepackage{caption}
\usepackage{subcaption}
\usepackage{hyperref}
\usepackage{url}

\captionsetup{font=small,labelfont=bf}

\title{What the Aleatoric Head Measures: A Residual Law for\\
Uncertainty Decomposition in Implicit Neural Representations}

% Hidden while the tmlr package is used without [accepted] or [preprint].
\author{\name Nishant Dave \email davenishu448@gmail.com \\
      \addr Independent Researcher}

\begin{document}
\maketitle

\begin{abstract}
\noindent
Predictive uncertainty is routinely split into an \emph{aleatoric} part, meant to capture
noise in the data, and an \emph{epistemic} part, meant to capture the model's own ignorance.
We ask what these two numbers measure in practice, on a task where every cause of
uncertainty is fixed by construction. Training one implicit neural representation over
$n$ images with a Gaussian negative-log-likelihood head and a five-member ensemble, across a
97-setting grid over data volume, capacity and noise, we find that the reported aleatoric
variance $A$ moves by $679\times$ while the true noise variance it estimates is held
constant over most of the grid, and lands within $1.5\times$ of that truth in only $13\%$ of
settings. One relation accounts for this: $A/\sigma^2$ is a power law in
$\mathrm{RMSE}_{\mathrm{train}}/\sigma$ with fitted exponent $1.86$--$1.92$
($R^2 = 0.93$--$0.96$), which is the optimum of the loss the head is trained on. The
aleatoric head is a meter for the training residual, and coincides with the noise only where
the fit happens to sit on the noise floor. Two pre-registered interventions establish this
causally: raising task difficulty at fixed capacity, data and noise moves $A$ by $6.4\times$
($\rho=+0.98$), and training the same model on the same data $27\times$ longer --- changing
nothing about the data, model or noise --- cuts the reported ``noise in the data'' by
$12.6\times$ ($\rho=-0.98$). In a companion analysis on a model whose posterior is available
in closed form, we show the epistemic term is no better off: an ensemble's spread converges
to $0.44\times$ the true epistemic variance when members differ by initialisation and
$6.72\times$ when they differ by resampling, never to $1.00\times$, while the two biases
partly cancel so that the \emph{undecomposed} total is accurate to $1.10\times$. We report
one pre-registered falsifier that failed. We recommend reporting total predictive variance,
and reporting $\mathrm{RMSE}_{\mathrm{train}}/\sigma$ beside any aleatoric number.
\end{abstract}

\section{Introduction}

A predictive uncertainty estimate is often decomposed into two parts
\citep{kendall2017,hullermeier2021}. The \emph{aleatoric} part is attributed to noise in the
observations and treated as irreducible; the \emph{epistemic} part is attributed to the
model's ignorance and expected to shrink with more data. The split is not cosmetic: it
prescribes opposite remedies. Aleatoric uncertainty says the measurements are the limit, so
buy better sensors; epistemic uncertainty says the model is the limit, so collect more data,
or label where the model is unsure.

Those prescriptions are only as good as the two numbers. This paper asks what they actually
measure, in a setting constructed so that the answer is checkable: a synthetic image family
with analytic ground truth, a noise level set by the experimenter, and a region withheld from
training in every image, so the amount of noise and the amount of missing information are
both known rather than inferred.

Our object of study is an implicit neural representation (INR) \citep{mildenhall2020,
sitzmann2020, tancik2020} with a heteroscedastic Gaussian output head \citep{nix1994,
kendall2017} and a deep ensemble \citep{lakshminarayanan2017}, which is the standard
construction for obtaining both terms at once. We sweep three knobs --- the number of images
the network must represent, its width, and the observation noise --- over 97 settings.

Our findings are:

\begin{enumerate}
\item \textbf{The aleatoric estimate does not track the noise} (\S\ref{sec:anomaly}). Over
the 48 settings where the noise level is identical by construction, $A$ still varies by
$490\times$. Across the grid, $\mathrm{corr}(\log\sigma^2, \log A) = +0.354$ while
$\mathrm{corr}(\log \mathrm{RMSE}_{\mathrm{train}}^2, \log A) = +0.933$.

\item \textbf{One law explains the grid} (\S\ref{sec:law}). The Gaussian NLL term is
minimised at the squared residual on the \emph{training} data. Accordingly,
$A/\sigma^2$ collapses onto a power law in $\mathrm{RMSE}_{\mathrm{train}}/\sigma$ with
exponent $\approx 2$ and $R^2 = 0.93$--$0.96$ in three independent reruns. Three regimes
follow: a model fitted below the noise floor under-reports noise ($0.26\times$), one above it
reports its own misfit as noise ($9.78\times$), and only a narrow band is correct.

\item \textbf{The relation is causal} (\S\ref{sec:causal}). Two pre-registered interventions
move the training residual while leaving the true noise untouched, and $A$ follows both,
landing on the same line the observational grid drew.

\item \textbf{The epistemic term is not a refuge} (\S\ref{sec:epistemic}). In a model with a
closed-form posterior, an ensemble's spread converges to a constant multiple of the true
epistemic variance --- $0.44\times$ or $6.72\times$ depending on how members are made to
differ --- and the multiple does not shrink with ensemble size.

\item \textbf{The sum survives better than its parts} (\S\ref{sec:total}), and the
decomposition's downstream value is limited (\S\ref{sec:cost}).
\end{enumerate}

We also report a pre-registered falsifier that failed, and the post-hoc analysis that
partially rescues it, at the same length as the successes (\S\ref{sec:causal}).

\section{Related work}

The aleatoric/epistemic distinction and its estimation are surveyed by
\citet{hullermeier2021}; \citet{kendall2017} give the formulation we study, in which a
network predicts a mean and a variance and an approximate posterior supplies the rest.
\citet{depeweg2018} develop the decomposition for decision-making, and
\citet{lakshminarayanan2017} introduce the deep ensemble whose spread is almost universally
read as the epistemic term. The variance head itself dates to \citet{nix1994}.

Several works report that the decomposition misbehaves. \citet{valdenegro2022} find
experimentally that the two terms interact, which they note ``violates assumptions on
aleatoric uncertainty'', and that aleatoric estimates are unreliable out of distribution.
\citet{mucsanyi2024} benchmark disentanglement across methods and find the specialised
uncertainty tasks substantially harder than predictive uncertainty, with estimators that
claim to separate the two failing to do so. \citet{bengs2022} prove a negative result on the
other side: the loss functions used to induce second-order predictors do not incentivise
faithful representation of epistemic uncertainty, so loss minimisation cannot be expected to
deliver it.

Closest to our mechanism is \citet{seitzer2022}, who show that Gaussian NLL training with
gradient-based optimisers can converge to poor-but-stable parameter estimates because the
gradient of the mean is scaled by the predicted variance, and propose $\beta$-NLL to
decouple them. Their concern is the \emph{optimisation} pathology --- the mean fit is damaged
en route. Ours is complementary and concerns the \emph{converged} estimate: even where
optimisation proceeds normally, the variance head's optimum is the training residual, so what
it reports is a property of the fit rather than of the data, and we give the quantitative law
and two causal interventions for that claim.

Relative to this literature our contribution is to replace the observation ``the terms
interact'' with a mechanism, a one-parameter law that collapses $679\times$ of variation, a
diagnostic that tells a practitioner when the aleatoric number is meaningful, and causal
evidence from interventions that hold the data and the noise fixed.

\section{Setup}
\label{sec:setup}

\paragraph{The world.} Images are defined analytically on $[0,1]^2$ as a sum of a linear
ramp, a Gaussian blob, a hard-edged disc and a sinusoidal grating, with per-image random
offsets scaled by a diversity parameter $D$ (default $D=1$). Because the images are
analytic, the clean value --- and therefore the true error --- is available at any
coordinate, not only on the pixel grid. Observations are made on a $32\times32$ grid with
additive Gaussian noise of standard deviation $\sigma$, specified as a percentage of the
images' own standard deviation. A central block, $|x-0.5|<0.12$ and $|y-0.5|<0.12$, is
excluded from training and calibration in every image; $20\%$ of the remaining pixels are
held out for conformal calibration \citep{vovk2005, angelopoulos2021}.

\paragraph{The model.} One INR $f(x, y, k) \to (\mu, \log v)$ maps a coordinate and an image
index to a mean and a log-variance. Inputs are mapped through $32$ random Fourier features
\citep{tancik2020, rahimi2007} at scale $1.5$ and concatenated with a learned $8$-dimensional
code per image; three hidden layers follow. Five independent networks are trained together as
one batched model. Training is AdamW, $10^{-3}$ weight decay, learning rate $2\times10^{-3}$
with cosine annealing, $1000$ steps, of which the first $25\%$ use squared error before
switching to the Gaussian NLL
\begin{equation}
\mathcal{L} = \tfrac12\left(\log v + \frac{(\mu - y)^2}{v}\right).
\label{eq:nll}
\end{equation}

\paragraph{The two quantities.} With $M=5$ members, writing $\mu_i$ and $v_i$ for member
$i$'s outputs, the standard decomposition \citep{kendall2017, valdenegro2022} is
\begin{equation}
\underbrace{\tfrac1M\textstyle\sum_i v_i}_{A \;\text{(aleatoric)}}
\;+\;
\underbrace{\mathrm{Var}_i[\mu_i]}_{E \;\text{(epistemic)}}
\;=\; \text{total predictive variance}.
\label{eq:decomp}
\end{equation}
We also record $\mathrm{RMSE}_{\mathrm{train}}$, the root-mean-square residual on the
training pixels; the true error at random off-grid query points; the rank correlation between
predicted standard deviation and true error; and raw and conformal coverage.

\paragraph{The grid.} Three knobs --- $n \in [1,100]$ images, width $\in [1,100]$ units,
noise $\in [1,100]\%$ --- are explored as three one-at-a-time sweeps plus three pairwise
heatmaps, giving 97 distinct settings, each with five networks. The whole grid was run three
times as the recording code was extended; the model, seeds and image family are identical
across the three, and we report all three throughout.

\paragraph{Pre-registration.} The hypotheses, thresholds and accept/reject rules of
\S\ref{sec:causal} were fixed in a dated document before those runs were executed. It is
included in the repository.

\section{The aleatoric estimate does not track the noise}
\label{sec:anomaly}

\begin{figure}[t]
\centering
\includegraphics[width=\textwidth]{figures/fig1_anomaly.png}
\caption{The quantity $A$ is meant to estimate (dashed) is constant by construction in two of
the three sweeps. $A$ (solid) is not. Log scales on both axes.}
\label{fig:anomaly}
\end{figure}

In the $n$ and width sweeps the noise level is never touched, so a faithful aleatoric
estimate should be flat. Figure~\ref{fig:anomaly} shows it is not, and
Table~\ref{tab:anomaly} quantifies the failure.

\begin{table}[h]
\centering
\small
\begin{tabular}{lccccc}
\toprule
run & $A$ spans & $\mathrm{corr}(\log\sigma^2,\log A)$ &
$\mathrm{corr}(\log \mathrm{RMSE}_{\mathrm{train}}^2, \log A)$ &
$A$ as est.\ of $\sigma^2$ & within $1.5\times$ \\
\midrule
1 & $637\times$ & $+0.391$ & $\mathbf{+0.958}$ & $6.74\times$ median & $18\%$ \\
2 & $679\times$ & $+0.354$ & $\mathbf{+0.933}$ & $6.28\times$ median & $13\%$ \\
3 & $679\times$ & $+0.354$ & $\mathbf{+0.933}$ & $6.28\times$ median & $13\%$ \\
\bottomrule
\end{tabular}
\caption{Across 97 settings, the quantity $A$ is defined to estimate explains almost none of
its variation, and the training residual explains nearly all of it. Over the 48 settings
where $\sigma^2$ is identical by construction, $A$ still ranges
$4.75\times10^{-5}$ to $2.33\times10^{-2}$, a factor of $490$.}
\label{tab:anomaly}
\end{table}

\section{One law}
\label{sec:law}

Minimising \eqref{eq:nll} in $v$ at fixed $\mu$ gives $v^\star = (\mu-y)^2$: the head's
optimum is the squared residual on the data it is trained on. If that is what it reports at
convergence, then whatever drives that residual --- too little capacity, too many images for
the capacity available, or too short a training run --- is reported as noise, and
\begin{equation}
\frac{A}{\sigma^2} \;\propto\; \left(\frac{\mathrm{RMSE}_{\mathrm{train}}}{\sigma}\right)^{2}.
\label{eq:law}
\end{equation}

\begin{figure}[t]
\centering
\includegraphics[width=0.78\textwidth]{figures/fig2_collapse.png}
\caption{All 97 grid settings (blue; marker shape indicates which sweep), plus the
intervention runs of \S\ref{sec:causal} (orange, green), collapse onto a single power law.
The dashed line is the exponent $2$ predicted by the NLL optimum.}
\label{fig:collapse}
\end{figure}

\begin{table}[h]
\centering
\small
\begin{tabular}{lccc}
\toprule
run & fitted exponent & 95\% CI & $R^2$ \\
\midrule
1 & $\mathbf{1.923}$ & $[1.84, 2.00]$ & $\mathbf{0.955}$ \\
2 & $\mathbf{1.864}$ & $[1.74, 1.98]$ & $\mathbf{0.932}$ \\
3 & $\mathbf{1.864}$ & $[1.74, 1.98]$ & $\mathbf{0.932}$ \\
\bottomrule
\end{tabular}
\caption{One free parameter collapses $679\times$ of variation in $A$ onto a single line, in
every rerun. Theory predicts exponent $2$. Intervals are percentile bootstrap over settings.}
\label{tab:law}
\end{table}

The fit holds (Table~\ref{tab:law}, Figure~\ref{fig:collapse}), which means the three sweeps
are one phenomenon rather than three. The ratio
$\mathrm{RMSE}_{\mathrm{train}}/\sigma$ --- where the fit sits relative to the noise floor ---
determines not just the magnitude but the sign of the error in $A$:

\begin{table}[h]
\centering
\small
\begin{tabular}{llcccl}
\toprule
regime & $\mathrm{RMSE}_{\mathrm{train}}/\sigma$ & settings & median $A/\sigma^2$ & median $\rho(\hat s, |e|)$ & where it occurs \\
\midrule
memorising   & $<0.8$        & 11 & $\mathbf{0.26\times}$ & $+0.17$ & few images; wide net, little data \\
matched      & $0.8$--$1.25$ & 13 & $\mathbf{0.75\times}$ & $+0.18$ & a narrow band only \\
underfitting & $>1.25$       & 73 & $\mathbf{9.78\times}$ & $+0.37$ & many images; narrow nets; short training \\
\bottomrule
\end{tabular}
\caption{A model fitted below the noise floor has absorbed noise into its weights and
under-reports it; one that has not reached the floor reports its own misfit as a property of
the world. $\rho(\hat s, |e|)$ is the rank correlation between predicted standard deviation
and true error (\S\ref{sec:cost}).}
\label{tab:regimes}
\end{table}

\section{Causal test}
\label{sec:causal}

The collapse is a correlation across settings that differ in several ways at once. Two
interventions, pre-registered before execution, hold the data, the model and the noise fixed
--- or as close to it as each allows --- and move only the residual. Both use the grid's own
baseline ($n=10$, width $50$, noise $10\%$) and three seeds.

\paragraph{I1: task difficulty.} The diversity parameter $D$ scales how much the images differ
from one another: a harder representation problem at identical capacity, data volume and noise
percentage. (Because the noise is specified as a percentage of image standard deviation,
$\sigma^2$ does move slightly; the law is stated in ratio form and we report both columns.)

\paragraph{I2: optimisation budget.} The number of training steps. This changes nothing about
the data, the model, or the noise --- only how far the fit is taken --- and is the cleaner
probe of the two.

\begin{table}[h]
\centering
\small
\begin{tabular}{cccccc@{\hskip 2em}cccccc}
\toprule
\multicolumn{6}{c}{\textbf{I1} --- task difficulty} & \multicolumn{6}{c}{\textbf{I2} --- optimisation budget} \\
\cmidrule(r){1-6}\cmidrule(l){7-12}
$D$ & $\frac{\mathrm{RMSE}_{\mathrm{tr}}}{\sigma}$ & $A$ & $\sigma^2$ & $\frac{A}{\sigma^2}$ & $E$ &
steps & $\frac{\mathrm{RMSE}_{\mathrm{tr}}}{\sigma}$ & $A$ & $\sigma^2$ & $\frac{A}{\sigma^2}$ & $E$ \\
\midrule
0   & 0.99 & 1.17e-4 & 1.57e-4 & $\mathbf{0.74}$ & 2.4e-5 &  150 & 4.83 & 3.03e-3 & 1.831e-4 & $\mathbf{16.56}$ & 3.08e-4 \\
0.5 & 2.17 & 2.45e-4 & 1.67e-4 & $\mathbf{1.47}$ & 1.5e-4 &  300 & 3.49 & 1.08e-3 & 1.831e-4 & $\mathbf{5.92}$  & 2.54e-4 \\
1   & 2.34 & 4.11e-4 & 1.83e-4 & $\mathbf{2.25}$ & 2.7e-4 &  600 & 2.64 & 5.46e-4 & 1.831e-4 & $\mathbf{2.98}$  & 2.82e-4 \\
2   & 2.16 & 5.49e-4 & 2.41e-4 & $\mathbf{2.27}$ & 2.6e-4 & 1000 & 2.34 & 4.11e-4 & 1.831e-4 & $\mathbf{2.25}$  & 2.66e-4 \\
3   & 1.94 & 6.75e-4 & 3.26e-4 & $\mathbf{2.07}$ & 2.7e-4 & 2000 & 2.14 & 3.34e-4 & 1.831e-4 & $\mathbf{1.82}$  & 2.64e-4 \\
    &      &         &         &                 &        & 4000 & 1.95 & 2.64e-4 & 1.831e-4 & $\mathbf{1.44}$  & 2.56e-4 \\
\bottomrule
\end{tabular}
\caption{Left: with every image identical ($D=0$) the fit lands on the noise floor and $A$ is
nearly correct; making the images differ raises the reported ``noise'' although the images'
noise is unchanged. Right: $\sigma^2$ is pinned at $1.831\times10^{-4}$ while training longer
cuts $A$ by $12.6\times$. Means over three seeds.}
\label{tab:interventions}
\end{table}

\paragraph{Pre-registered outcomes.}
\textbf{H1} (accepted): $\rho(D, A) = +0.982$, range $6.4\times$ (threshold: $\rho>0.7$ and
range $>2\times$).
\textbf{H2} (accepted): $\rho(\text{steps}, A) = -0.981$, range $12.6\times$ (threshold:
$\rho<-0.7$). The quantity of noise a model reports in its data is a function of how long it
was trained.
\textbf{H3} (accepted): pooling the 33 intervention runs, the fitted exponent is $1.939$
$[1.53, 2.48]$ with $R^2 = 0.882$, against $1.864$ from the observational grid and $2$ from
theory --- the interventions land on the line the grid drew (Figure~\ref{fig:collapse}).
\textbf{H4} (accepted): $\mathrm{RMSE}_{\mathrm{train}}^2$ predicts $A$ to a median factor of
$2.22$ versus $2.74$ for the error at held-out test points, consistent with the head reading
training residuals rather than generalisation error. We note against ourselves that on these
33 runs $\sigma^2$ alone scores $2.14$, better than both; the interventions sit at one point
of the grid, so $\sigma^2$ barely varies there and a constant fits adequately. The grid is
where that comparison is informative (Table~\ref{tab:anomaly}).

\paragraph{H5, the falsifier, failed.} We pre-registered that $E$ must \emph{not} obey a
comparable law, since the mechanism is specific to the variance head; the criterion was
$|\text{slope}|<1.0$ or $R^2<0.5$. Pooled over both interventions, $E$ gives slope $1.382$
with $R^2=0.638$, so \textbf{H5 is rejected as written}.

A post-hoc split, which we label as such, shows what the pooled test conflated. I1 changes
the images themselves, which genuinely changes how much there is for members to disagree
about; only I2 isolates the head. Splitting the same regression:

\begin{center}
\small
\begin{tabular}{lcccc}
\toprule
runs & $A$: slope, $R^2$ & $E$: slope, $R^2$ & $A$ range & $E$ range \\
\midrule
I2 (budget only)     & $2.64$, $0.987$ & $0.14$, $0.294$ & $\mathbf{12.6\times}$ & $\mathbf{1.4\times}$ \\
I1 (task difficulty) & $1.21$, $0.827$ & $2.43$, $0.970$ & $6.4\times$ & $12.3\times$ \\
\bottomrule
\end{tabular}
\end{center}

Under I2 alone --- train $27\times$ longer, change nothing else --- $A$ moves $12.6\times$
and $E$ moves $1.4\times$, the asymmetry the mechanism predicts. This is a post-hoc rescue of
a pre-registered test that failed, and is weaker evidence than a pass would have been. Our
reading is that I1 was never a clean probe of the aleatoric head alone, and that writing the
pooled criterion down was our error rather than a discovery about $E$.

\section{The epistemic term is not a refuge}
\label{sec:epistemic}

If $A$ is compromised, the natural response is to rely on $E$. To test that we need the true
epistemic variance, which the INR does not provide. We therefore replace the trained layers
above the INR's fixed random Fourier features with a Bayesian linear head. The resulting
random-feature Bayesian linear regression is a close analogue of the same model --- same
input featurisation, same smoothness prior --- with posterior covariance
$\Sigma = (\Phi^\top W \Phi + \alpha I)^{-1}$, $W=\mathrm{diag}(1/\sigma_i^2)$, so the exact
epistemic variance at a point is $\phi^\top\Sigma\phi$ and the exact aleatoric variance is the
known noise map. We use a $64\times64$ field with noise rising $0.02 \to 0.12$ left to right,
$8\%$ of pixels observed, an unobserved block, a band sampled at $0.25\times$ density, $512$
features, and 30 independent draws. Hypotheses were pre-registered; three of four failed.

\begin{figure}[t]
\centering
\includegraphics[width=\textwidth]{figures/exact_convergence.png}
\caption{Left: only an ensemble of exact posterior samples converges at the theoretical rate.
Right: each construction converges to a different multiple of the true epistemic variance.}
\label{fig:exact}
\end{figure}

\begin{table}[h]
\centering
\small
\begin{tabular}{lcc}
\toprule
members differ by & estimate / truth at $M=50$ & error decays as \\
\midrule
exact posterior samples & $\mathbf{1.00\times}$ & $M^{-0.457}$ (theory $-0.5$) \\
random initialisation \emph{(a deep ensemble)} & $\mathbf{0.44\times}$ & $M^{-0.102}$ --- a floor \\
bootstrap resampling & $\mathbf{6.72\times}$ & a floor \\
\bottomrule
\end{tabular}
\caption{An ensemble's spread estimates the variance of whatever distinguishes its members.
That equals the posterior's epistemic variance only when the members are posterior samples.}
\label{tab:exact}
\end{table}

Table~\ref{tab:exact} and Figure~\ref{fig:exact} give the result: with members that really
are posterior samples the estimator behaves exactly as theory says, and with the two standard
constructions it converges to the wrong quantity, missing in opposite directions so that no
fixed correction is available to a practitioner who does not already know the answer. We
pre-registered that the finite-$M$ error would decay as $(M-1)^{-1/2}$; it does not, because
most of the gap is not sampling error.

The entanglement is also definitional rather than estimation-induced. Doubling the
observation noise --- a purely aleatoric intervention --- multiplies the \emph{exact}
epistemic variance by $3.63$ $[3.61, 3.65]$, against the ceiling of $4.00$ that this model's
algebra predicts. The INR shows the same effect where it can be measured: across its noise
sweep $E$ rises $13.3\times$ from $1\%$ to $100\%$ noise, and a doubling at the high end
($32\%\to66\%$) multiplies it by $2.69\times$.

\section{The sum survives better than its parts}
\label{sec:total}

\begin{center}
\small
\begin{tabular}{lc}
\toprule
quantity (exact-posterior model, $M=50$) & estimate / truth \\
\midrule
epistemic term & $0.44\times$ \\
aleatoric term & $1.43\times$ \\
\textbf{total predictive variance} & $\mathbf{1.10\times}$ \\
\bottomrule
\end{tabular}
\end{center}

The two errors run in opposite directions and largely cancel, leaving the undecomposed total
the best-estimated quantity available. The INR grid agrees from the other side: the rank
correlation between $A$ and $E$ across points within a setting has median $+0.457$ and
reaches $-1.000$, so the split is not merely noisy --- its sign is unstable. Splitting a
quantity estimable to $10\%$ into two parts estimable to a factor of two to three is a poor
trade unless the split buys something.

\section{What it costs}
\label{sec:cost}

\begin{figure}[t]
\centering
\includegraphics[width=\textwidth]{figures/fig3_consequence.png}
\caption{Left: training longer reduces the reported ``noise in the data'' by $12.6\times$
while the true noise (dashed) is fixed and the epistemic term is flat. Right: the rank
correlation between predicted uncertainty and true error collapses as noise grows.}
\label{fig:cost}
\end{figure}

Let $\rho(\hat s, |e|)$ be the per-setting rank correlation between predicted total standard
deviation and true absolute error: does the uncertainty point at the mistakes? Across the
grid it runs $+0.03$ to $+0.85$ with median $+0.31$. It is \emph{never negative} in any of the
96 settings where it is defined, so the uncertainty degrades to uninformative rather than
inverting. Its correlation with the noise knob is $-0.64$: along the noise sweep it falls from
$+0.52$ at $1\%$ noise to $+0.03$ at $100\%$ --- that is, it is least useful exactly where a
user would most want a warning.

Calibration does not repair this, and repairs less than it appears to. In the exact-posterior
model, split-conformal scaling \citep{vovk2005, angelopoulos2021} delivers its marginal
guarantee precisely --- global coverage $0.900$, inside $[0.88,0.92]$ in 27 of 30 draws ---
while coverage inside the unobserved block swings between $0.24$ and $0.97$ across
statistically identical draws. Tracing that swing: the true extrapolation error inside the
block varies $7.4\times$ across draws while the predicted uncertainty there varies $1.9\times$,
and the two are uncorrelated ($r=-0.01$, against $r=-0.75$ for the actual error). In the
region the model never observed, the uncertainty estimate carries almost no information about
the error it purports to describe, and an exact Bayesian posterior does not fix this: a
posterior's variance in an unobserved region is set by the prior and the sampling geometry,
not by what the unseen signal does there.

\section{Recommendations}

\begin{enumerate}
\item \textbf{Report the total.} It is accurate to $\approx 1.10\times$ where the parts are
$0.44\times$ and $1.43\times$ with unstable sign (\S\ref{sec:total}).
\item \textbf{Report $\mathrm{RMSE}_{\mathrm{train}}/\sigma$ beside any aleatoric number, or
report no aleatoric number.} It is free to compute and is the single statistic that says
whether $A$ is meaningful: below $0.8$ the model has absorbed the noise, above $1.25$ it is
reporting its own misfit (\S\ref{sec:law}).
\item \textbf{Never compare $A$ across models, training budgets or dataset sizes.} A
$27\times$ change in training budget moves it $12.6\times$ with the data untouched
(\S\ref{sec:causal}).
\item \textbf{Do not add ensemble members to improve the epistemic estimate.} The error has a
floor that $M$ does not reduce (\S\ref{sec:epistemic}).
\item \textbf{Read conformal coverage as the marginal guarantee it is.} It says nothing about
any particular region, least of all one without data (\S\ref{sec:cost}).
\end{enumerate}

\section{Limitations}

The data are synthetic throughout, and the noise is Gaussian and heteroscedastic exactly as
the models assume; real misspecification would likely make the picture worse rather than
better. Two model families are studied and one of them is an analogue: random-feature
Bayesian linear regression resembles the INR but its exactness comes from linearity in the
weights, and a trained INR's ensembles fail to be posterior samples in a different way than
studied here. The causal evidence is for $A$; for $E$ we have the exact-posterior analysis
and, in the INR itself, only the post-hoc asymmetry of \S\ref{sec:causal}. One pre-registered
falsifier failed. Finally, we found and documented but did not fix a defect in the image
generator: for $D \geq 3.33$ its blob-width term can go negative and produce non-finite
pixels; all grids reported here were produced at $D=1$ and the intervention sweep stops at
$D=3$.

\section*{Reproducibility}

Code, pre-registrations, completed grids and the scripts that regenerate every number and
figure in this paper are available at an anonymised repository for review:
\url{https://anonymous.4open.science/r/ANON}. It will be de-anonymised on acceptance. The
grid runs on a laptop GPU; the interventions and the exact-posterior study run on CPU in
about 21 and 15 minutes respectively.

\begin{thebibliography}{99}
\small

\bibitem[Angelopoulos and Bates(2021)]{angelopoulos2021}
A.~N. Angelopoulos and S.~Bates.
\newblock A gentle introduction to conformal prediction and distribution-free uncertainty
quantification.
\newblock \emph{arXiv:2107.07511}, 2021.

\bibitem[Bengs et al.(2022)]{bengs2022}
V.~Bengs, E.~H\"ullermeier, and W.~Waegeman.
\newblock Pitfalls of epistemic uncertainty quantification through loss minimisation.
\newblock In \emph{Advances in Neural Information Processing Systems 35 (NeurIPS)}, 2022.
arXiv:2203.06102.

\bibitem[Depeweg et al.(2018)]{depeweg2018}
S.~Depeweg, J.~M. Hern\'andez-Lobato, F.~Doshi-Velez, and S.~Udluft.
\newblock Decomposition of uncertainty in Bayesian deep learning for efficient and risk-sensitive
learning.
\newblock In \emph{International Conference on Machine Learning (ICML)}, 2018.
arXiv:1710.07283.

\bibitem[H\"ullermeier and Waegeman(2021)]{hullermeier2021}
E.~H\"ullermeier and W.~Waegeman.
\newblock Aleatoric and epistemic uncertainty in machine learning: An introduction to concepts
and methods.
\newblock \emph{Machine Learning}, 110(3):457--506, 2021.

\bibitem[Kendall and Gal(2017)]{kendall2017}
A.~Kendall and Y.~Gal.
\newblock What uncertainties do we need in Bayesian deep learning for computer vision?
\newblock In \emph{Advances in Neural Information Processing Systems 30 (NIPS)}, pages
5574--5584, 2017. arXiv:1703.04977.

\bibitem[Lakshminarayanan et al.(2017)]{lakshminarayanan2017}
B.~Lakshminarayanan, A.~Pritzel, and C.~Blundell.
\newblock Simple and scalable predictive uncertainty estimation using deep ensembles.
\newblock In \emph{Advances in Neural Information Processing Systems 30 (NIPS)}, 2017.
arXiv:1612.01474.

\bibitem[Mildenhall et al.(2020)]{mildenhall2020}
B.~Mildenhall, P.~P. Srinivasan, M.~Tancik, J.~T. Barron, R.~Ramamoorthi, and R.~Ng.
\newblock NeRF: Representing scenes as neural radiance fields for view synthesis.
\newblock In \emph{European Conference on Computer Vision (ECCV)}, 2020. arXiv:2003.08934.

\bibitem[Mucs\'anyi et al.(2024)]{mucsanyi2024}
B.~Mucs\'anyi, M.~Kirchhof, and S.~J. Oh.
\newblock Benchmarking uncertainty disentanglement: Specialized uncertainties for specialized
tasks.
\newblock \emph{arXiv:2402.19460}, 2024.

\bibitem[Nix and Weigend(1994)]{nix1994}
D.~A. Nix and A.~S. Weigend.
\newblock Estimating the mean and variance of the target probability distribution.
\newblock In \emph{Proceedings of the IEEE International Conference on Neural Networks
(ICNN)}, pages 55--60, 1994.

\bibitem[Rahimi and Recht(2007)]{rahimi2007}
A.~Rahimi and B.~Recht.
\newblock Random features for large-scale kernel machines.
\newblock In \emph{Advances in Neural Information Processing Systems 20 (NIPS)}, 2007.

\bibitem[Seitzer et al.(2022)]{seitzer2022}
M.~Seitzer, A.~Tavakoli, D.~Antic, and G.~Martius.
\newblock On the pitfalls of heteroscedastic uncertainty estimation with probabilistic neural
networks.
\newblock In \emph{International Conference on Learning Representations (ICLR)}, 2022.
arXiv:2203.09168.

\bibitem[Sitzmann et al.(2020)]{sitzmann2020}
V.~Sitzmann, J.~N.~P. Martel, A.~W. Bergman, D.~B. Lindell, and G.~Wetzstein.
\newblock Implicit neural representations with periodic activation functions.
\newblock In \emph{Advances in Neural Information Processing Systems 33 (NeurIPS)}, 2020.
arXiv:2006.09661.

\bibitem[Tancik et al.(2020)]{tancik2020}
M.~Tancik, P.~P. Srinivasan, B.~Mildenhall, S.~Fridovich-Keil, N.~Raghavan, U.~Singhal,
R.~Ramamoorthi, J.~T. Barron, and R.~Ng.
\newblock Fourier features let networks learn high frequency functions in low dimensional
domains.
\newblock In \emph{Advances in Neural Information Processing Systems 33 (NeurIPS)}, 2020.
arXiv:2006.10739.

\bibitem[Valdenegro-Toro and Saromo Mori(2022)]{valdenegro2022}
M.~Valdenegro-Toro and D.~Saromo Mori.
\newblock A deeper look into aleatoric and epistemic uncertainty disentanglement.
\newblock In \emph{IEEE/CVF Conference on Computer Vision and Pattern Recognition Workshops
(CVPRW)}, pages 1509--1517, 2022. arXiv:2204.09308.

\bibitem[Vovk et al.(2005)]{vovk2005}
V.~Vovk, A.~Gammerman, and G.~Shafer.
\newblock \emph{Algorithmic Learning in a Random World}.
\newblock Springer, 2005.

\end{thebibliography}

\end{document}
